\chapter{拓扑连通性}
\label{chap:topological-connectedness}

区间介值定理所使用的整体性质是：区间不能分成两个互不接触的非空开部分。把这一性质从实直线中抽出，便得到连通性。它由开闭集刻画，因而适用于一般拓扑空间；连续映射保持连通性，所以介值现象也随之推广。

道路连通要求任意两点之间有一条连续曲线。欧氏空间中的许多常见集合使“连通”与“道路连通”看起来几乎相同，但一般子空间中二者会分开。本章给出这个差别的经典例子，并说明局部道路连通性怎样重新使两种分支一致。

\section{闭开分离与连通分支}

\begin{definition}[分离与连通]
拓扑空间 \(X\) 的一个分离，是一对不相交非空开集 \(U,V\)，满足
\[
 X=U\cup V.
\]
不存在分离的空间称为连通空间。子集 \(A\subseteq X\) 称为连通集，如果它带子空间拓扑时连通。
\end{definition}

在分离中，\(U=X\setminus V\)、\(V=X\setminus U\)，所以两部分同时开闭。于是 \(X\) 连通，当且仅当它的开闭子集只有 \(\varnothing\) 和 \(X\)。

\begin{theorem}[连通集的粘合与闭包]
\label{thm:connected-union-closure}
\begin{enumerate}[label=\textup{(\roman*)}]
 \item 若连通子集族 \((C_i)_{i\in I}\) 有一个公共点，则其并集
 \(\bigcup_{i\in I}C_i\) 连通；
 \item 连通集 \(C\) 的闭包 \(\overline C\) 连通；
 \item 更一般地，若每个 \(C_i\) 都与一个固定连通集 \(C\) 相交，则
 \(C\cup\bigcup_iC_i\) 连通。
\end{enumerate}
\end{theorem}

\begin{proof}
证明(iii)即可。反设并集 \(E\) 有分离 \(U,V\)。固定连通集 \(C\) 必完全落在其中一侧，不妨 \(C\subseteq U\)。每个 \(C_i\) 与 \(C\) 相交，故与 \(U\) 相交；由 \(C_i\) 连通，它不能同时与 \(V\) 相交，所以 \(C_i\subseteq U\)。于是 \(V\) 为空，矛盾。(i)是取 \(C\) 为公共单点的特例。

对(ii)，若 \(\overline C=U\cup V\) 是分离，则 \(C\) 完全落在一侧，不妨 \(C\subseteq U\)。此时相对闭集 \(\overline C\setminus U=V\) 与 \(C\) 不交，但 \(C\) 在 \(\overline C\) 中稠密，故 \(V=\varnothing\)，矛盾。
\end{proof}

\begin{theorem}[递减紧连通集的交]
\label{thm:nested-compact-connected-intersection}
设 \(X\) 为 Hausdorff 空间，
\[
 K_1\supseteq K_2\supseteq\cdots
\]
是一列非空紧连通子集。则
\[
 K=\bigcap_{n\in\Npos}K_n
\]
非空且连通。
\end{theorem}

\begin{proof}
各 \(K_n\) 是紧集 \(K_1\) 中的闭集，并具有有限交性质，所以
\(K\ne\varnothing\)。

若 \(K\) 不连通，存在 \(K\) 中不相交非空开闭集 \(A,B\) 覆盖 \(K\)。
它们作为紧 Hausdorff 空间 \(K\) 的闭子集是紧集，因而也是 \(X\) 中的闭集。紧 Hausdorff 空间 \(K_1\) 正规，所以可在 \(K_1\) 中取不相交开集 \(U,V\)，分别包含 \(A,B\)。

若每个 \(K_n\) 都有点落在 \(K_1\setminus(U\cup V)\)，则闭集
\[
 K_n\setminus(U\cup V)
\]
在 \(K_1\) 中非空、紧且递减；有限交性质会给出一点属于
\(K\setminus(U\cup V)\)，与 \(K=A\cup B\subseteq U\cup V\) 矛盾。因此某个 \(K_N\subseteq U\cup V\)。另一方面，\(K_N\) 同时与
\(A\subseteq U\) 和 \(B\subseteq V\) 相交，于是
\[
 K_N=(K_N\cap U)\cup(K_N\cap V)
\]
是一个分离，违背 \(K_N\) 连通。
\end{proof}

\begin{definition}[连通分支]
给定 \(x\in X\)，所有含 \(x\) 的连通子集之并称为 \(x\) 的连通分支，记为 \(C(x)\)。
\end{definition}

由粘合定理，\(C(x)\) 连通；它包含所有含 \(x\) 的连通集，所以是含 \(x\) 的最大连通子集。

\begin{proposition}[连通分支的结构]
\label{prop:connected-components}
关系
\[
 x\sim y\quad\Longleftrightarrow\quad
 \text{存在连通子集同时含 }x,y
\]
是等价关系，其等价类正是连通分支。每个连通分支闭，但不一定开。
\end{proposition}

\begin{proof}
每个单点集连通，所以关系自反；把同一连通集中的两点交换次序，得到对称性。若 \(x,y\) 同属连通集 \(A\)，而 \(y,z\) 同属连通集 \(B\)，则 \(A\cap B\ne\varnothing\)，故 \(A\cup B\) 连通，得到 \(x\sim z\)。所以它是等价关系，等价类按定义是最大连通集。

连通分支 \(C\) 的闭包仍连通，并含 \(C\)。由极大性 \(\overline C=C\)，故 \(C\) 闭。它未必开：\(\Q\) 的每个连通分支都是单点，而 \(\Q\) 没有孤立点。
\end{proof}

最后一句将在习题中证明。它提醒我们：一个空间可以由闭的连通分支分解，却不能把这些分支当作开块逐个研究。

\section{道路分支与局部道路连通}

\begin{definition}[道路与道路连通]
从 \(x\) 到 \(y\) 的道路是连续映射
\[
 \gamma:[0,1]\to X,\qquad \gamma(0)=x,\quad\gamma(1)=y.
\]
若任意两点之间都有道路，则称 \(X\) 道路连通。含一点 \(x\) 的最大道路连通子集称为 \(x\) 的道路分支。
\end{definition}

常值道路、反向参数化和分段重参数化分别给出道路关系的自反、对称与传递性，因此道路分支也是等价类。

\begin{proposition}
\label{prop:path-connected-implies-connected}
道路连通空间连通。因而每个道路分支包含在某个连通分支中。
\end{proposition}

\begin{proof}
若 \(X=U\cup V\) 是分离并取 \(x\in U,y\in V\)，从 \(x\) 到 \(y\) 的道路 \(\gamma\) 会把连通区间 \([0,1]\) 连续映到一个同时交 \(U,V\) 的连通集，矛盾。第二句由道路分支本身连通得到。
\end{proof}

\begin{definition}[局部连通与局部道路连通]
若每个点都有由连通开邻域组成的邻域基，则称空间局部连通。若每个点都有由道路连通开邻域组成的邻域基，则称空间局部道路连通。
\end{definition}

“每点有一个连通邻域”不足以构成局部连通；要求的是在任意给定邻域内还能缩小到连通开邻域。

\begin{theorem}[局部道路连通空间的分支]
\label{thm:locally-path-connected-components}
若 \(X\) 局部道路连通，则每个道路分支开，并且道路分支与连通分支相同。更一般地，\(X\) 局部道路连通，当且仅当 \(X\) 的每个开子集的道路分支都是开集。
\end{theorem}

\begin{proof}
设 \(P\) 是 \(X\) 的道路分支，\(x\in P\)。取道路连通开邻域 \(U\ni x\)。\(U\) 中每一点可先沿 \(U\) 内道路到 \(x\)，再沿 \(P\) 中道路到固定基点，故 \(U\subseteq P\)。所以 \(P\) 开。道路分支把 \(X\) 分成互不相交开集。固定一个道路分支 \(P\)，它的补集是其余开道路分支的并，也开；所以 \(P\) 是开闭集。任一连通集若与 \(P\) 相交，就不能再与 \(X\setminus P\) 相交。因此每个连通分支只含一个道路分支，两者相同。

开子空间继承局部道路连通性，故正向对每个开子集成立。反之，给 \(x\in U\) 且 \(U\) 开；\(x\) 在 \(U\) 中的道路分支按假设开，是含 \(x\) 且包含于 \(U\) 的道路连通开邻域。这些分支构成所需邻域基。
\end{proof}

把“道路”删去，同一证明给出：空间局部连通，当且仅当每个开子集的连通分支开。局部道路连通蕴含局部连通，逆命题一般不成立。

\section{连通但非道路连通的空间}

\subsection{正弦曲线的闭包与道路分支}

令
\[
 G=\{(x,\sin(1/x)):0<x\le1\}\subseteq\R^2,
\qquad
 S=G\cup(\{0\}\times[-1,1]).
\]
集合 \(S\) 称为拓扑学家的正弦曲线。

\begin{theorem}[拓扑学家的正弦曲线]
\label{thm:topologists-sine-curve}
集合 \(S\) 连通但不道路连通。它恰有两个道路分支：曲线 \(G\) 与竖直线段 \(\{0\}\times[-1,1]\)。
\end{theorem}

\begin{proof}
\theoremref{thm:topologists-sine-curve-not-path-connected}已经在度量空间中
证明 \(S=\overline G\) 连通，并利用振荡序列排除了从竖直线段到 \(G\) 的
道路。那里还证明了 \(G\) 与竖直线段各自道路连通，故它们恰为两个道路分支。
\end{proof}

由于 \(S=\overline G\) 且 \(S\) 在 \(\R^2\) 中闭、有界，它还是一个紧致
连通空间。曲线 \(G\) 是道路分支，却不是 \(S\) 中的闭集，因为它的闭包为
整个 \(S\)。这与连通分支形成对照：连通分支总是闭的，道路分支则未必闭。
因此，把空间按道路分支取商时，所得商空间甚至可能不是 \(T_1\) 空间；本例中
对应于 \(G\) 的单点原像不是闭集。这个现象来自竖直线段附近无穷多次积聚，
而不是由任一局部距离估计造成。

\section{连续像、商空间与局部连通性}

\subsection{连续像与商空间}

\begin{theorem}[连续像保持连通]
\label{thm:continuous-image-connected}
若 \(X\) 连通、\(f:X\to Y\) 连续，则 \(f(X)\) 连通。
\end{theorem}

\begin{proof}
若 \(f(X)\) 有分离 \(U,V\)，则 \(f^{-1}(U),f^{-1}(V)\) 是 \(X\) 的不相交非空相对开集，并覆盖 \(X\)，矛盾。
\end{proof}

当 \(Y=\R\) 时，连通子集恰为区间，因而上一定理包含介值定理：连通空间上的连续实值函数，其值域是区间。它没有使用定义域中的距离或次序。

\begin{corollary}[商映射与连通性]
\label{cor:quotient-connected}
连通空间的任意商空间连通。更一般地，若 \(q:X\to Y\) 连续且满射，则 \(X\) 连通蕴含 \(Y\) 连通。
\end{corollary}

\begin{proof}
\(Y=q(X)\)，直接应用连续像定理。
\end{proof}

\subsection{任意乘积}

\begin{theorem}[乘积连通性]
\label{thm:product-connected}
任意一族非空连通空间的积在积拓扑下连通。
\end{theorem}

\begin{proof}
选定基点 \(a=(a_i)\in\prod_iX_i\)。对有限集 \(F\subseteq I\)，令
\[
 X_F=\{x:x_i=a_i\text{ 对所有 }i\notin F\}.
\]
\(X_F\) 同胚于有限多个连通空间之积。先看两个因子：固定
\(y_0\in Y\)，水平集 \(X\times\{y_0\}\) 连通；对每个 \(x\in X\)，
竖直集 \(\{x\}\times Y\) 连通并与该水平集相交。把这些集合粘合便得到
\(X\times Y\) 连通，有限积再由归纳得到。因此每个 \(X_F\) 连通。
所有 \(X_F\) 都含 \(a\)，所以
\[
 X_{\mathrm{fin}}=\bigcup_{F\subseteq I,\ F\text{ 有限}}X_F
\]
连通。

\(X_{\mathrm{fin}}\) 在积空间中稠密，因为任一基本开集只限制有限多个坐标，可在这些坐标内选点，其余坐标取 \(a_i\)。故整个积是连通集 \(X_{\mathrm{fin}}\) 的闭包，因而连通。选取基点使用了从各非空因子同时选点的选择原则。
\end{proof}

\subsection{局部连通性的保持}

局部连通性在连续像下不能无条件保持；开映射提供了一个常用充分条件。

\begin{proposition}[开映射保持局部连通]
\label{prop:open-map-local-connectedness}
设 \(q:X\to Y\) 连续、满射且为开映射。若 \(X\) 局部连通（或局部道路连通），则 \(Y\) 也具有相应性质。
\end{proposition}

\begin{proof}
给 \(y\in Y\) 的开邻域 \(V\)，取 \(x\in q^{-1}(y)\)。原像
\(q^{-1}(V)\) 开；在其中取含 \(x\) 的连通开邻域 \(U\)（或道路连通开邻域）。则 \(q(U)\) 连通（或道路连通），又因 \(q\) 开而是 \(y\) 的开邻域，并包含于 \(V\)。
\end{proof}

商映射不一定开，所以商空间虽自动保持整体连通性，却未必保持局部连通性。使用识别构造空间时，应分别检查商映射的连续性、满射性和开映射性质。

\section*{习题}
\addcontentsline{toc}{section}{习题}

\subsection*{A组　连通集与分支}
\begin{enumerate}[label=\textbf{\thechapter.\arabic*.}]
 \exerciseitem{connected-clopen-functions}
 证明 \(X\) 连通，当且仅当每个连续映射 \(f:X\to\{0,1\}\)（目标带离散拓扑）都是常值。
 \exerciseitem{connected-intersecting-family}
 设 \((C_i)\) 均连通，并且存在 \(i_0\) 使每个 \(C_i\cap C_{i_0}\ne\varnothing\)。证明 \(\bigcup_iC_i\) 连通；再给出“只要求任意相邻两个相交”但并集不连通的错误表述，并说明需要怎样组织指标才正确。
 \exerciseitem{connected-subsets-real-line}
 证明 \(\R\) 的非空连通子集恰好是各种区间，包括单点、开闭半开区间和无界区间。
 \exerciseitem{components-rational-line}
 证明 \(\Q\) 的每个连通分支都是单点，但单点在 \(\Q\) 中不开。
 \exerciseitem{component-map}
 若 \(f:X\to Y\) 连续，证明每个连通分支的像包含在 \(Y\) 的某个连通分支中。说明一般不能把“包含”改成“等于”。
\end{enumerate}

\subsection*{B组　道路与局部性质}
\begin{enumerate}[label=\textbf{\thechapter.\arabic*.},resume]
 \exerciseitem{path-concatenation-continuity}
 写出两条首尾相接道路的拼接公式，并用粘合引理证明其连续性。
 \exerciseitem{path-components-open-criterion}
 证明空间局部道路连通，当且仅当每个开子集的道路分支都开。
 \exerciseitem{locally-connected-components-open}
 证明空间局部连通，当且仅当每个开子集的连通分支都开。
 \exerciseitem{convex-star-path-connected}
 证明赋范线性空间中的凸集道路连通；更一般地，证明星形集道路连通。
 \exerciseitem{deleted-plane-paths}
 证明 \(\R^n\setminus\{0\}\) 在 \(n\ge2\) 时道路连通，而 \(\R\setminus\{0\}\) 有两个连通分支。
 \exerciseitem{sine-curve-closure-detail}
 逐项证明拓扑学家的正弦曲线 \(S\) 恰为 \(G\) 的闭包。再沿用
 \theoremref{thm:topologists-sine-curve-not-path-connected} 中的
 \(u,h,t_0\) 记号，说明时刻 \(s_n,r_n\) 为什么可以选取，以及为什么
 \(s_n,r_n\to t_0\)。
\end{enumerate}

\subsection*{C组　构造与保持性}
\begin{enumerate}[label=\textbf{\thechapter.\arabic*.},resume]
 \exerciseitem{connected-product-finite}
 不使用正文的任意乘积定理，先证明两个连通空间的积连通，再归纳到有限积。
 \exerciseitem{connected-product-dense-finite-support}
 补足任意乘积证明中 \(X_{\mathrm{fin}}\) 稠密的基论证，并说明箱拓扑下为什么同一证明失效。
 \exerciseitem{quotient-interval-circle-connected}
 由商映射证明圆周 \(S^1\) 连通，并证明有限多个圆周在一点粘合所得空间连通。
 \exerciseitem{open-quotient-local-path}
 若 \(q:X\to Y\) 是开商映射且 \(X\) 局部道路连通，证明 \(Y\) 局部道路连通。构造一个商映射不是开映射的简单例子。
 \optionalexerciseitem{continuous-image-components-not-surjective}
 构造连续满射 \(f:X\to Y\) 和 \(X\) 的一个连通分支 \(C\)，使 \(f(C)\) 严格包含于 \(Y\) 中的某个连通分支。
\end{enumerate}
